% GENERATED FILE. Edit canonical lessons, then run npm run generate:books.
% canonical-source-hash: fnv1a64:628c7fe51f14b060
% canonical-lessons: TA-C456-muccu, TA-C456-perumuccu, TA-C456-nimmati, TA-C456-arutal, TA-C456-tani

\chapter{Breath and Relief}
\label{ch:ta-breath-and-relief}
\hlchaptermodality{456}

\begin{chapteropening}
\textbf{By the end of this chapter:} \emph{I can talk about breathing, sighing and feeling relieved.}

Complete the last lesson of chapter 456: I can talk about breathing, sighing and feeling relieved.
\end{chapteropening}

\section[\texorpdfstring{\ta{மூச்சு} (mūccu)}{மூச்சு (mūccu)}]{\ta{மூச்சு} (mūccu) — breath}

\label{lesson:TA-C456-muccu}



\begin{quote}
Before the new one: say the Tamil for to bemoan, then the Tamil for to suffer.
\end{quote}

\subsection*{You'll want to know: \ta{மூச்சு}}
\textbf{\ta{மூச்சு}} — \emph{mūccu} — \textquotedblleft{}breath\textquotedblright{}.

\begin{quote}
\ta{ஓடி} \ta{வந்ததால்} \ta{எனக்கு} \ta{மூச்சு} \ta{வாங்கியது}.
\end{quote}

\begin{quote}
\emph{ōṭi vantatāl eṉakku mūccu vāṅkiyatu.}
\end{quote}

\begin{quote}
I was out of breath from running.
\end{quote}

\ta{மூச்சு} \ta{வாங்கு}, to pant; \ta{மூச்சு} \ta{விடு}, to breathe out.

\begin{cousinweb}
\textbf{\ta{மூச்சு}} (\emph{mūccu}) is native Dravidian (DEDR 4886); the entry joins it with verbs of smelling and sniffing.
\end{cousinweb}

\begin{grammarlens}[title={What you've built}]
One more word for describing what you lived and felt.
\end{grammarlens}

\subsection*{Your turn}
\begin{itemize}
\raggedright
  \item \emph{Say it:} \emph{mūccu}
  \item \emph{Say it:} \emph{mūccu}, once more
  \item \emph{Say it:} \emph{ōṭi vantatāl eṉakku mūccu vāṅkiyatu.} — I was out of breath from running.
  \item \emph{Recall:} say the Tamil for to bemoan, then the Tamil for to suffer, then say \emph{mūccu} again
\end{itemize}

\begin{tcolorbox}[breakable,colback=teal!4,colframe=teal!35!black,title={Before you move on}]
What does \ta{மூச்சு} mean? (\textquotedblleft{}Breath\textquotedblright{}.) Say it once more.
\end{tcolorbox}
\section[\texorpdfstring{\ta{பெருமூச்சு} (perumūccu)}{பெருமூச்சு (perumūccu)}]{\ta{பெருமூச்சு} (perumūccu) — a sigh}

\label{lesson:TA-C456-perumuccu}



\begin{quote}
Before the new one: say the Tamil for to suffer, then the Tamil for breath.
\end{quote}

\subsection*{You'll want to know: \ta{பெருமூச்சு}}
\textbf{\ta{பெருமூச்சு}} — \emph{perumūccu} — \textquotedblleft{}a sigh\textquotedblright{}.

\begin{quote}
\ta{வேலை} \ta{முடிந்ததும்} \ta{அவள்} \ta{ஒரு} \ta{பெருமூச்சு} \ta{விட்டாள்}.
\end{quote}

\begin{quote}
\emph{vēlai muṭintatum avaḷ oru perumūccu viṭṭāḷ.}
\end{quote}

\begin{quote}
When the work was over she heaved a sigh.
\end{quote}

\ta{பெரு} \textquotedblleft{}big\textquotedblright{} + \ta{மூச்சு} \textquotedblleft{}breath\textquotedblright{}; \ta{பெருமூச்சு} \ta{விடு}, to sigh.

\begin{cousinweb}
\textbf{\ta{பெருமூச்சு}} (\emph{perumūccu}) is a native compound of \ta{பெரு} \emph{peru}, \textquotedblleft{}great\textquotedblright{} (DEDR 4411), and \ta{மூச்சு} \emph{mūccu}, \textquotedblleft{}breath\textquotedblright{} (DEDR 4886): a big breath.
\end{cousinweb}

\begin{grammarlens}[title={What you've built}]
One more word for describing what you lived and felt.
\end{grammarlens}

\subsection*{Your turn}
\begin{itemize}
\raggedright
  \item \emph{Say it:} \emph{perumūccu}
  \item \emph{Say it:} \emph{perumūccu}, once more
  \item \emph{Say it:} \emph{vēlai muṭintatum avaḷ oru perumūccu viṭṭāḷ.} — When the work was over she heaved a sigh.
  \item \emph{Recall:} say the Tamil for to suffer, then the Tamil for breath, then say \emph{perumūccu} again
\end{itemize}

\begin{tcolorbox}[breakable,colback=teal!4,colframe=teal!35!black,title={Before you move on}]
What does \ta{பெருமூச்சு} mean? (\textquotedblleft{}A sigh\textquotedblright{}.) Say it once more.
\end{tcolorbox}
\section[\texorpdfstring{\ta{நிம்மதி} (nimmati)}{நிம்மதி (nimmati)}]{\ta{நிம்மதி} (nimmati) — relief; peace of mind}

\label{lesson:TA-C456-nimmati}



\begin{quote}
Before the new one: say the Tamil for breath, then the Tamil for a sigh.
\end{quote}

\subsection*{You'll want to know: \ta{நிம்மதி}}
\textbf{\ta{நிம்மதி}} — \emph{nimmati} — \textquotedblleft{}relief; peace of mind\textquotedblright{}.

\begin{quote}
\ta{மகன்} \ta{வீட்டுக்கு} \ta{வந்த} \ta{பிறகுதான்} \ta{அம்மாவுக்கு} \ta{நிம்மதி} \ta{வந்தது}.
\end{quote}

\begin{quote}
\emph{makaṉ vīṭṭukku vanta piṟakutāṉ ammāvukku nimmati vantatu.}
\end{quote}

\begin{quote}
Only after her son came home was the mother at peace.
\end{quote}

\ta{நிம்மதியாக}, with peace of mind: \ta{நிம்மதியாகத்} \ta{தூங்கினேன்}, I slept in peace.

\begin{cousinweb}
\textbf{\ta{நிம்மதி}} (\emph{nimmati}) is borrowed from Telugu, a sister Dravidian language: DEDR 3602 lists it, \textquotedblleft{}tranquillity\textquotedblright{}, in an entry of words for desire and affection.
\end{cousinweb}

\begin{grammarlens}[title={What you've built}]
One more word for describing what you lived and felt.
\end{grammarlens}

\subsection*{Your turn}
\begin{itemize}
\raggedright
  \item \emph{Say it:} \emph{nimmati}
  \item \emph{Say it:} \emph{nimmati}, once more
  \item \emph{Say it:} \emph{makaṉ vīṭṭukku vanta piṟakutāṉ ammāvukku nimmati vantatu.} — Only after her son came home was the mother at peace.
  \item \emph{Recall:} say the Tamil for breath, then the Tamil for a sigh, then say \emph{nimmati} again
\end{itemize}

\begin{tcolorbox}[breakable,colback=teal!4,colframe=teal!35!black,title={Before you move on}]
What does \ta{நிம்மதி} mean? (\textquotedblleft{}Relief; peace of mind\textquotedblright{}.) Say it once more.
\end{tcolorbox}
\section[\texorpdfstring{\ta{ஆறுதல்} (āṟutal)}{ஆறுதல் (āṟutal)}]{\ta{ஆறுதல்} (āṟutal) — consolation; comfort}

\label{lesson:TA-C456-arutal}



\begin{quote}
Before the new one: say the Tamil for a sigh, then the Tamil for relief.
\end{quote}

\subsection*{You'll want to know: \ta{ஆறுதல்}}
\textbf{\ta{ஆறுதல்}} — \emph{āṟutal} — \textquotedblleft{}consolation; comfort\textquotedblright{}.

\begin{quote}
\ta{உன்} \ta{வார்த்தைகள்} \ta{எனக்கு} \ta{ஆறுதலாக} \ta{இருந்தன}.
\end{quote}

\begin{quote}
\emph{uṉ vārttaikaḷ eṉakku āṟutalāka iruntaṉa.}
\end{quote}

\begin{quote}
Your words were a comfort to me.
\end{quote}

From \ta{ஆறு} \textquotedblleft{}to cool, to heal\textquotedblright{} with the noun ending -\ta{தல்}. \ta{ஆறுதல்} \ta{சொல்}, to console.

\begin{cousinweb}
\textbf{\ta{ஆறுதல்}} (\emph{āṟutal}) is native Dravidian, a noun from \ta{ஆறு} \emph{āṟu}, \textquotedblleft{}to cool, to be appeased\textquotedblright{} (DEDR 404): comfort is a cooling of pain.
\end{cousinweb}

\begin{grammarlens}[title={What you've built}]
One more word for describing what you lived and felt.
\end{grammarlens}

\subsection*{Your turn}
\begin{itemize}
\raggedright
  \item \emph{Say it:} \emph{āṟutal}
  \item \emph{Say it:} \emph{āṟutal}, once more
  \item \emph{Say it:} \emph{uṉ vārttaikaḷ eṉakku āṟutalāka iruntaṉa.} — Your words were a comfort to me.
  \item \emph{Recall:} say the Tamil for a sigh, then the Tamil for relief, then say \emph{āṟutal} again
\end{itemize}

\begin{tcolorbox}[breakable,colback=teal!4,colframe=teal!35!black,title={Before you move on}]
What does \ta{ஆறுதல்} mean? (\textquotedblleft{}Consolation; comfort\textquotedblright{}.) Say it once more.
\end{tcolorbox}
\section[\texorpdfstring{\ta{தணி} (taṇi)}{தணி (taṇi)}]{\ta{தணி} (taṇi) — to subside; to abate}

\label{lesson:TA-C456-tani}



\begin{quote}
Before the new one: say the Tamil for relief, then the Tamil for consolation.
\end{quote}

\subsection*{You'll want to know: \ta{தணி}}
\textbf{\ta{தணி}} — \emph{taṇi} — \textquotedblleft{}to subside; to abate\textquotedblright{}.

\begin{quote}
\ta{மாத்திரை} \ta{சாப்பிட்ட} \ta{பிறகு} \ta{காய்ச்சல்} \ta{தணிந்தது}.
\end{quote}

\begin{quote}
\emph{māttirai cāppiṭṭa piṟaku kāyccal taṇintatu.}
\end{quote}

\begin{quote}
After the tablet the fever subsided.
\end{quote}

Past (it …): \ta{தணிந்தது}, \emph{taṇintatu}.

Of fever, heat or anger: \ta{கோபம்} \ta{தணிந்தது}, the anger cooled. Said of the thing, so its past is \ta{தணிந்தது}.

\begin{cousinweb}
\textbf{\ta{தணி}} (\emph{taṇi}) is native Dravidian (DEDR 3045), in the family of \ta{தண்} \emph{taṇ}, \textquotedblleft{}cool\textquotedblright{}.
\end{cousinweb}

\begin{grammarlens}[title={What you've built}]
One more word for describing what you lived and felt.
\end{grammarlens}

\subsection*{Your turn}
\begin{itemize}
\raggedright
  \item \emph{Say it:} \emph{taṇi}
  \item \emph{Say it:} \emph{taṇi}, once more
  \item \emph{Say it:} \emph{māttirai cāppiṭṭa piṟaku kāyccal taṇintatu.} — After the tablet the fever subsided.
  \item \emph{Recall:} say the Tamil for relief, then the Tamil for consolation, then say \emph{taṇi} again
\end{itemize}

\begin{tcolorbox}[breakable,colback=teal!4,colframe=teal!35!black,title={Before you move on}]
What does \ta{தணி} mean? (\textquotedblleft{}To subside; to abate\textquotedblright{}.) Say it once more.
\end{tcolorbox}
